\documentclass[11pt]{article}

\usepackage[margin=1in]{geometry}
\usepackage{amsmath,amssymb}
\usepackage{booktabs}
\usepackage{graphicx}
\usepackage{url}
\usepackage[hidelinks]{hyperref}
\usepackage{microtype}

\title{Lost in One Index: Measuring How Language and Modality Gaps\\
Compound in English--Indonesian Multimodal Scholarly Search}

\author{Shaban Lubanga\\
Department of Information Systems\\
Institut Teknologi Sepuluh Nopember, Surabaya, Indonesia\\
\texttt{6026251019@student.its.ac.id}}

\date{}

\begin{document}
\maketitle

\begin{abstract}
Scholarly knowledge in Indonesia is split across languages: research papers
are published in English, while lectures and theses are produced in
Indonesian. Retrieval systems that embed all sources into a single vector
index promise to bridge this divide, but their cross-lingual behavior is
rarely measured. We build a bilingual multimodal test collection --- ten
English papers (2{,}720 text chunks, 67 figures), eight English and eight
Indonesian lecture videos (2{,}101 transcript segments), and 15
translation-paired queries --- and measure how retrieval changes when the
same question is asked in English versus Indonesian. Under an English-only
CLIP encoder, retrieval silos by language: paired queries return nearly
disjoint top-10 lists (mean Jaccard overlap 0.019), Indonesian queries reach
English papers in only 3 of 15 cases, and the language and modality gaps
compound --- Indonesian-query-to-figure similarity is the lowest-scoring
group in the collection (mean 0.498 vs.\ 0.812 for same-language
transcripts). A multilingual CLIP text encoder closes most of the
consistency gap (Jaccard 0.474) but introduces a register bias that floods
results with lecture transcripts. Per-group score calibration restores
coverage --- figures return to up to 7 of 15 queries and every Indonesian
query reaches English papers --- but halves relevance for Indonesian queries
(nDCG@10 $0.120 \rightarrow 0.051$), exposing a fairness--precision
tradeoff. A BM25 baseline outperforms every dense configuration in both
languages, aided by code-switched English terminology in Indonesian
lectures. We release the collection, relevance judgments, and code.
\end{abstract}

\section{Introduction}

A student at an Indonesian university studying deep learning lives between
two languages. The primary literature --- arXiv papers, conference
proceedings --- is in English. The explanations that make that literature
accessible --- university lectures on YouTube, campus theses, national
journals --- are largely in Bahasa Indonesia. Writing a literature review
means searching both worlds, yet search tools handle them separately: an
Indonesian question does not find English papers, and an English query does
not surface the Indonesian lecture that best explains the concept.

Multimodal retrieval systems that embed text, figures, and lecture
transcripts into one shared vector space appear to offer a way out: if all
content lives in the same space, one query should retrieve across languages
and modalities alike. Prior work on such systems, including our own earlier
prototype~\cite{synthesizer}, has documented a \emph{modality gap}: in a
shared CLIP index, figures score systematically below text and are almost
never retrieved~\cite{liang2022mind}. This paper asks what happens when a
\emph{language} dimension is added, using English--Indonesian as the test
pair. Indonesian is spoken by more than 270 million people, yet
Indonesian-language evaluation resources remain scarce~\cite{nusacrowd,
indonlu}, and to our knowledge no prior work measures cross-lingual behavior
in multimodal scholarly search.

We make four contributions:
\begin{enumerate}
  \item \textbf{A bilingual multimodal test collection.} Ten English
  deep-learning papers, eight English and eight Indonesian lecture videos,
  15 translation-paired queries (authored and verified by a native
  speaker), and 999 pooled relevance judgments, all released publicly.
  \item \textbf{A measurement of the language gap and its interaction with
  the modality gap.} Under English-only CLIP, translation-paired queries
  return nearly disjoint results (Jaccard 0.019), and the two gaps stack:
  Indonesian queries against English figures form the lowest-similarity
  group in the collection.
  \item \textbf{An encoder comparison.} A multilingual CLIP text encoder
  recovers most cross-lingual consistency (Jaccard 0.474) but compresses
  similarity scores and biases retrieval toward spoken-register content,
  displacing paper text.
  \item \textbf{A calibration experiment revealing a fairness--precision
  tradeoff.} Z-normalizing scores within each (modality, language) group
  restores coverage across groups but reduces relevance for Indonesian
  queries by half, showing that balanced exposure and precision are
  competing objectives in a single shared index.
\end{enumerate}

\section{Related Work}

\textbf{Cross-lingual and multilingual retrieval.} Cross-language
information retrieval has classically relied on translation~\cite{clir};
multilingual sentence encoders such as LaBSE~\cite{labse} and
knowledge-distilled multilingual CLIP~\cite{mclip} enable direct
cross-lingual dense retrieval by aligning languages in one embedding space.
Benchmarks such as MIRACL~\cite{miracl} evaluate multilingual text
retrieval, and Indonesian NLP resources are expanding through efforts like
IndoNLU~\cite{indonlu} and NusaCrowd~\cite{nusacrowd}. Multimodal
\emph{and} multilingual scholarly retrieval, however, remains unmeasured;
our collection targets exactly this intersection.

\textbf{The modality gap.} Liang et al.~\cite{liang2022mind} showed that
contrastively trained multimodal encoders place image and text embeddings
in systematically different regions of the space. In retrieval over
academic content, this manifests as figure starvation: in our earlier
system, figures were never retrieved by unified top-$k$ search over a
shared CLIP index~\cite{synthesizer}. The present work measures how this
gap interacts with a language gap and evaluates one mitigation.

\textbf{Multimodal document retrieval.} ColPali~\cite{colpali} retrieves
document pages as images with vision-language embeddings, and
M3DocRAG~\cite{m3docrag} performs retrieval-augmented generation over
multimodal document collections. These systems target visually rich English
documents; our setting adds spoken lectures and a second language.

\textbf{LLM-based relevance judgment.} Recent work shows large language
models can produce relevance judgments that correlate with human assessors
at useful levels~\cite{thomas2024, umbrela}. We adopt this methodology,
with guidelines, per-decision conventions, and known biases documented
alongside the released judgments.

\section{Test Collection and System}

\subsection{Corpus}
The corpus covers deep learning broadly, chosen so that Indonesian lecture
material exists for every topic:

\begin{itemize}
  \item \textbf{Papers (English):} ten arXiv papers ---
  word2vec~\cite{word2vec}, VGG~\cite{vgg}, ResNet~\cite{resnet},
  Transformer~\cite{transformer}, BERT~\cite{bert}, Mamba~\cite{mamba},
  Jamba~\cite{jamba}, Mamba-360~\cite{mamba360}, Mamba-2~\cite{mamba2},
  and an empirical study of Mamba-based LMs~\cite{empiricalmamba} ---
  yielding 2{,}720 text chunks and 67 embedded raster figures.
  \item \textbf{Lectures:} eight English videos (including 3Blue1Brown,
  Stanford CS231N/CS224N, and Michigan lectures) and eight Indonesian
  videos from university and community channels (Machine Learning
  Indonesia, Telkom University course recordings, and others), yielding
  1{,}146 English and 955 Indonesian transcript segments. Indonesian
  transcripts are auto-generated captions; transcription noise is part of
  the measured condition.
  \item \textbf{Queries:} 15 study questions in translation-paired form
  (English and Indonesian), spanning architectures, training methods, and
  model comparisons. The Indonesian phrasings were written and then
  corrected against academic register norms by the author, a native
  speaker.
\end{itemize}

Text and transcripts are chunked to at most 50 words to respect the CLIP
text encoder's 77-token input limit. Figures carry captions parsed from PDF
text or generated by LLaVA-7B~\cite{llava}, and are embedded as the average
of image and caption embeddings. All items carry modality and language
metadata. The index contains 4{,}888 vectors of dimension 512, searched
exactly with FAISS~\cite{faiss} inner product over unit-normalized vectors.

\subsection{Retrieval configurations}
We compare, at $k{=}10$ over identical corpora:
\begin{itemize}
  \item \textbf{CLIP}: English CLIP ViT-B/32 for text and images.
  \item \textbf{mCLIP}: a multilingual text encoder distilled into the CLIP
  image space~\cite{mclip}, paired with the same image encoder.
  \item \textbf{BM25}: lexical retrieval~\cite{bm25} over every textual
  representation (chunks, captions, transcripts), both languages mixed.
  \item \textbf{Calibrated variants}: for each dense encoder, similarity
  scores are z-normalized per query \emph{within each (modality, language)
  group} before ranking, so that no group is advantaged by the scale of its
  score distribution.
\end{itemize}

\subsection{Measures}
For every configuration and query language we report: the modality and
content-language distribution of retrieved items; per-query coverage
(queries returning $\geq$1 figure, $\geq$1 item of each content language);
\emph{EN$\leftrightarrow$ID consistency} --- for each translation-paired
query, the Jaccard overlap between the English query's and the Indonesian
query's result sets, averaged over pairs (a language-fair system approaches
1); and graded effectiveness (P@10, nDCG@10) against pooled binary
relevance judgments. Judgments cover all 999 items retrieved by any
configuration and were produced by a large language model following
written guidelines (strict standard; term mention $\neq$ relevance; items
judged on meaning regardless of language; figures judged by inspecting the
image). Guidelines, per-convention notes, and known biases are released
with the data; the single-judge design is a limitation
(Section~\ref{sec:limits}).

\section{Results}

\subsection{English-only CLIP silos retrieval by language}
\label{sec:silo}

Table~\ref{tab:dist} shows the modality and language distribution of
retrieved items. Under CLIP, English queries retrieve almost exclusively
English content (148/150 items) and Indonesian queries almost exclusively
Indonesian content (134/150); only 3 of 15 Indonesian queries retrieve any
English paper content at all. The consistency between translation-paired
queries is near zero (Jaccard 0.019, overlap@10 0.033): \emph{the same
question, asked in the other language, returns an almost entirely different
result list}. Figures are never retrieved by either query language,
replicating the modality gap on a corpus $2.3\times$ larger than in our
earlier study~\cite{synthesizer}.

\begin{table}[t]
\centering
\caption{Retrieved items by modality and content language (15 queries
$\times$ top-10). Coverage counts queries returning $\geq$1 figure /
$\geq$1 English-content item / $\geq$1 Indonesian-content item.}
\label{tab:dist}
\small
\begin{tabular}{llrrrrrl}
\toprule
Config & Query lang. & Text & Fig. & Trans. & EN & ID & Coverage (fig/EN/ID)\\
\midrule
CLIP  & EN & 74 & 0 & 76  & 148 & 2   & 0/15 \;\; 15/15 \;\; 2/15\\
CLIP  & ID & 11 & 0 & 139 & 16  & 134 & 0/15 \;\; 3/15 \;\; 15/15\\
mCLIP & EN & 13 & 0 & 137 & 22  & 128 & 0/15 \;\; 8/15 \;\; 15/15\\
mCLIP & ID & 5  & 0 & 145 & 8   & 142 & 0/15 \;\; 5/15 \;\; 15/15\\
BM25  & EN & 85 & 4 & 61  & 146 & 4   & 2/15 \;\; 15/15 \;\; 1/15\\
BM25  & ID & 18 & 2 & 130 & 25  & 125 & 1/15 \;\; 7/15 \;\; 15/15\\
\bottomrule
\end{tabular}
\end{table}

\subsection{The gaps compound}
\label{sec:gaps}

Table~\ref{tab:gap} reports the mean cosine similarity between queries and
\emph{all} indexed items, grouped by (modality, content language). Two
structures are visible. First, the \emph{language gap}: under CLIP,
Indonesian queries score same-language transcripts at 0.812 but English
paper text at only 0.649 --- a same-language bonus of $+0.16$ that explains
the siloing in Section~\ref{sec:silo}. Second, the \emph{modality gap}:
figures score far below all text groups for both query languages. The two
effects stack: the Indonesian-query$\to$English-figure group (0.498) is the
lowest cell in the table, and its \emph{maximum} (0.622) barely reaches the
\emph{mean} of the competing same-language transcript group. Content that
is both cross-language and cross-modality is doubly disadvantaged.

\begin{table}[t]
\centering
\caption{Query--item cosine similarity by (modality, content language),
mean $\pm$ std over 15 queries $\times$ all items in the group.}
\label{tab:gap}
\small
\begin{tabular}{lcccc}
\toprule
 & \multicolumn{2}{c}{CLIP} & \multicolumn{2}{c}{mCLIP}\\
Group ($n$) & EN queries & ID queries & EN queries & ID queries\\
\midrule
Text EN (2720)       & $.714 \pm .075$ & $.649 \pm .072$ & $.835 \pm .033$ & $.843 \pm .035$\\
Figure EN (67)       & $.556 \pm .043$ & $.498 \pm .041$ & $.615 \pm .034$ & $.617 \pm .033$\\
Transcript EN (1146) & $.761 \pm .049$ & $.703 \pm .050$ & $.860 \pm .024$ & $.877 \pm .026$\\
Transcript ID (955)  & $.718 \pm .051$ & $.812 \pm .056$ & $.880 \pm .023$ & $.899 \pm .027$\\
\bottomrule
\end{tabular}
\end{table}

\subsection{Multilingual encoding closes the language gap but opens a
register gap}

Under mCLIP the per-group means become nearly identical for English and
Indonesian queries (Table~\ref{tab:gap}, right), and consistency between
paired queries rises from 0.019 to 0.474 Jaccard (overlap@10 0.593) --- the
encoder substitution alone removes most of the language siloing. But the
distributions also compress (standard deviations halve) and reorder:
transcripts of \emph{both} languages now outscore paper text, and retrieval
floods with lecture segments (137--145 of 150 items;
Table~\ref{tab:dist}), displacing the papers that hold the most
authoritative content. Figures remain at zero. The language bias is traded
for a register bias, and the underlying lesson generalizes: \emph{raw
similarity scores are not comparable across content groups, and which group
wins is an artifact of the encoder's training distribution.}

\subsection{Calibration restores coverage but costs precision}

Z-normalizing scores within each (modality, language) group forces the
groups into comparable ranges. Coverage responds immediately: figures
appear for up to 7 of 15 queries (mCLIP-calibrated: 5/15 for Indonesian
queries), every Indonesian query retrieves English paper content
(3/15 $\rightarrow$ 15/15 under CLIP), and the transcript flood recedes
(mCLIP, ID queries: 145 $\rightarrow$ 48 transcript items). Paired-query
consistency also peaks (Jaccard 0.505, overlap@10 0.640, mCLIP-calibrated).

Relevance tells the other half of the story (Table~\ref{tab:rel}).
Calibration leaves English-query effectiveness roughly unchanged (CLIP
nDCG@10 $.238 \rightarrow .259$; mCLIP $.227 \rightarrow .218$) but
\emph{halves} it for Indonesian queries (CLIP $.120 \rightarrow .051$;
mCLIP $.155 \rightarrow .092$). Forcing starved groups into the top-10
promotes items that are on-group but off-topic: with only 67 figures and a
strict relevance standard, a guaranteed figure slot is usually filled by an
irrelevant figure. Balanced exposure and precision are competing
objectives; a single shared index satisfies at most one at a time.

\begin{table}[t]
\centering
\caption{Effectiveness against pooled relevance judgments
(999 items, 127 relevant).}
\label{tab:rel}
\small
\begin{tabular}{lcccc}
\toprule
 & \multicolumn{2}{c}{P@10} & \multicolumn{2}{c}{nDCG@10}\\
Configuration & EN & ID & EN & ID\\
\midrule
BM25              & \textbf{.253} & \textbf{.180} & \textbf{.328} & \textbf{.241}\\
CLIP raw          & .193 & .113 & .238 & .120\\
CLIP calibrated   & .193 & .047 & .259 & .051\\
mCLIP raw         & .160 & .113 & .227 & .155\\
mCLIP calibrated  & .160 & .067 & .218 & .092\\
\bottomrule
\end{tabular}
\end{table}

\subsection{BM25 outperforms all dense configurations}

BM25 achieves the best effectiveness in both query languages
(Table~\ref{tab:rel}) while retrieving \emph{some} cross-modal content (via
figure captions) and even some cross-lingual content: 7 of 15 Indonesian
queries reach English material lexically. The bridge is code-switching ---
Indonesian technical lectures routinely embed English terms
(\emph{attention}, \emph{transformer}, \emph{layer}), giving lexical
matching shared vocabulary across the language boundary. On a
domain-specific corpus of this size, dense retrieval's semantic
generalization does not yet pay for its calibration problems.

\section{Discussion}

Our measurements support three practical conclusions. First,
\emph{single-index dense retrieval is unsafe across languages and
modalities without calibration}: which content gets retrieved is governed
by encoder-induced score offsets (language, register, modality), not only
by relevance. Second, \emph{each mitigation trades one bias for another}:
multilingual encoding removes the language offset but introduces a register
offset; per-group calibration equalizes exposure but promotes irrelevant
items from sparse groups. Systems that need both fairness and precision
should retrieve per group and merge with a relevance-aware policy ---
learned fusion, or lexical--dense hybrids that exploit code-switching ---
rather than rank raw or naively calibrated scores in one list. Third, for
low-resource-language deployments, \emph{BM25 remains a strong default},
and evaluations of multilingual dense retrievers should always include it.

\section{Limitations}
\label{sec:limits}

The collection is small (one topic domain, 4{,}888 items, 15 query pairs)
and covers one language pair; effect sizes are large and consistent, but
generalization to other languages and domains is untested. Relevance was
judged by a single LLM assessor; although the methodology follows recent
practice~\cite{thomas2024,umbrela} and all guidelines, conventions, and
judgments are released for audit, human verification --- even of a sample
--- would strengthen the conclusions. Indonesian transcripts are
auto-generated captions, so ASR noise and caption quality are confounded
with language effects (we treat this as part of the realistic condition,
since it is what any deployed system would index). Finally, snippet
truncation during judging may undercount the relevance of long text chunks
relative to transcripts.

\section{Conclusion}

We presented the first measurement, to our knowledge, of how language and
modality gaps interact in multimodal scholarly retrieval, using a new
English--Indonesian test collection. An English-only encoder silos
retrieval by language; a multilingual encoder fixes alignment but not score
comparability; calibration fixes exposure but not relevance; and a lexical
baseline quietly outperforms them all. The collection, judgments, and code
are released to support work on the missing piece: retrieval that is
simultaneously language-fair, modality-fair, and precise.

\section*{Acknowledgments}

The retrieval pipeline extends the Multimodal Research Synthesizer, built
jointly with Hippolyte Marc-Antoine Catteau-Verniers, who kindly agreed to
this solo extension. Experiment tooling, analysis scripts, and drafting
were assisted by an AI system (Claude, Anthropic); the author designed the
study, verified all results, and takes full responsibility for the content.

\begin{thebibliography}{99}

\bibitem{synthesizer} S.~Lubanga and H.~M.-A. Catteau-Verniers,
``Multimodal Research Synthesizer,'' software and technical report, 2026.
\url{https://github.com/shaban2/Multimodal-Research-Synthesizer}

\bibitem{liang2022mind} W.~Liang, Y.~Zhang, Y.~Kwon, S.~Yeung, and J.~Zou,
``Mind the Gap: Understanding the Modality Gap in Multi-modal Contrastive
Representation Learning,'' in \emph{Proc.\ NeurIPS}, 2022.

\bibitem{clir} J.-Y. Nie, \emph{Cross-Language Information Retrieval}.
Morgan \& Claypool, 2010.

\bibitem{labse} F.~Feng, Y.~Yang, D.~Cer, N.~Arivazhagan, and W.~Wang,
``Language-agnostic BERT Sentence Embedding,'' in \emph{Proc.\ ACL}, 2022.

\bibitem{mclip} N.~Reimers and I.~Gurevych, ``Making Monolingual Sentence
Embeddings Multilingual using Knowledge Distillation,'' in \emph{Proc.\
EMNLP}, 2020.

\bibitem{miracl} X.~Zhang et al., ``MIRACL: A Multilingual Retrieval
Dataset Covering 18 Diverse Languages,'' \emph{TACL}, vol.~11, 2023.

\bibitem{indonlu} B.~Wilie et al., ``IndoNLU: Benchmark and Resources for
Evaluating Indonesian Natural Language Understanding,'' in \emph{Proc.\
AACL-IJCNLP}, 2020.

\bibitem{nusacrowd} S.~Cahyawijaya et al., ``NusaCrowd: Open Source
Initiative for Indonesian NLP Resources,'' in \emph{Findings of ACL}, 2023.

\bibitem{colpali} M.~Faysse et al., ``ColPali: Efficient Document Retrieval
with Vision Language Models,'' in \emph{Proc.\ ICLR}, 2025.

\bibitem{m3docrag} J.~Cho, D.~Mahata, O.~\.Irsoy, Y.~He, and M.~Bansal,
``M3DocRAG: Multi-modal Retrieval is What You Need for Multi-page
Multi-document Understanding,'' arXiv:2411.04952, 2024.

\bibitem{thomas2024} P.~Thomas, S.~Spielman, N.~Craswell, and B.~Mitra,
``Large Language Models Can Accurately Predict Searcher Preferences,'' in
\emph{Proc.\ SIGIR}, 2024.

\bibitem{umbrela} S.~Upadhyay, R.~Pradeep, N.~Thakur, N.~Craswell, and
J.~Lin, ``UMBRELA: UMbrela is the (Open-Source Reproduction of the) Bing
RELevance Assessor,'' arXiv:2406.06519, 2024.

\bibitem{word2vec} T.~Mikolov, K.~Chen, G.~Corrado, and J.~Dean,
``Efficient Estimation of Word Representations in Vector Space,''
arXiv:1301.3781, 2013.

\bibitem{vgg} K.~Simonyan and A.~Zisserman, ``Very Deep Convolutional
Networks for Large-Scale Image Recognition,'' in \emph{Proc.\ ICLR}, 2015.

\bibitem{resnet} K.~He, X.~Zhang, S.~Ren, and J.~Sun, ``Deep Residual
Learning for Image Recognition,'' in \emph{Proc.\ CVPR}, 2016.

\bibitem{transformer} A.~Vaswani et al., ``Attention Is All You Need,'' in
\emph{Proc.\ NeurIPS}, 2017.

\bibitem{bert} J.~Devlin, M.-W. Chang, K.~Lee, and K.~Toutanova, ``BERT:
Pre-training of Deep Bidirectional Transformers for Language
Understanding,'' in \emph{Proc.\ NAACL}, 2019.

\bibitem{mamba} A.~Gu and T.~Dao, ``Mamba: Linear-Time Sequence Modeling
with Selective State Spaces,'' arXiv:2312.00752, 2023.

\bibitem{jamba} O.~Lieber et al., ``Jamba: A Hybrid Transformer-Mamba
Language Model,'' arXiv:2403.19887, 2024.

\bibitem{mamba360} B.~N. Patro et al., ``Mamba-360: Survey of State Space
Models as Transformer Alternative for Long Sequence Modelling,''
arXiv:2404.16112, 2024.

\bibitem{mamba2} T.~Dao and A.~Gu, ``Transformers are SSMs: Generalized
Models and Efficient Algorithms Through Structured State Space Duality,''
in \emph{Proc.\ ICML}, 2024.

\bibitem{empiricalmamba} R.~Waleffe et al., ``An Empirical Study of
Mamba-based Language Models,'' arXiv:2406.07887, 2024.

\bibitem{llava} H.~Liu, C.~Li, Q.~Wu, and Y.~J. Lee, ``Visual Instruction
Tuning,'' in \emph{Proc.\ NeurIPS}, 2023.

\bibitem{faiss} J.~Johnson, M.~Douze, and H.~J\'egou, ``Billion-scale
Similarity Search with GPUs,'' \emph{IEEE Trans.\ Big Data}, 2019.

\bibitem{bm25} S.~Robertson and H.~Zaragoza, ``The Probabilistic Relevance
Framework: BM25 and Beyond,'' \emph{Foundations and Trends in Information
Retrieval}, vol.~3, no.~4, 2009.

\end{thebibliography}

\end{document}
